\documentclass{article}
\usepackage{amsmath, amssymb}
\usepackage{bbm}
\usepackage{cancel}
\usepackage{tikz}
\usepackage{tikz-cd}
\usepackage{mdframed}
\usepackage{hyperref}
\usepackage{tikz}

\begin{document}

\section{Lecture 20: Plan for Class}
Last class
\begin{enumerate}
    \item Ergodic theory
\end{enumerate}

\section{Recap of where we were with Ergodic theory}

We have some measure-preserving system $\mu(A)=\mu(T^{-1}A)$ where $T:\, M\to M$ where $M$ is a measure space. Remember that a Hilbert space $H=L^2(M)$ where the norm is a square-integrable functions $\langle f,g\rangle = \int f(x)g(x)\, d\mu(x)$.

Now we can define an operator $U:H\to H$ where 
\[
Uf(x)=f(Tx)
\]
and we showed that $U$ is an isometry
\[
\left|\left| Uf \right|\right| = \left|\left| f \right|\right|
\]
We now define the space of invariant functions on $H$ such that 
\[
H^{\text{inv}} = \left\{ f\in H:\, Uf=f \right\}
\]
Then we defined an orthogonal projection onto $H^{\text{inv}}$ as $P$.

von Neumann's ergodic theorem was essentially a description of $P$ as 
\[
\lim_{n\to\infty}\,\frac{1}{n}\left(f+ Uf + U^2 f +\ldots + U^{n-1}f\right) = Pf
\]

\section{Finishing the proof of von Neumann's theory}

To start, let 
\[
G=\left\{ Ug - g :\, g\in H \right\}
\]
We first showed that $H^{\text{inv}}= G^\bot$. Then it follows that 
\[
H = H^{\text{inv}}\oplus G
\]
If follows that 
\begin{itemize}
    \item If $h\in H^{\text{inv}}$, then $\frac{1}{n}\left(h+ Uh + U^2 h +\ldots + U^{n-1}h\right) = h = Ph$
    \item If $g\in G$, then $\frac{1}{n}\left(g+ Ug + U^2 g +\ldots + U^{n-1}g\right) \to 0$ as $n\to \infty$
\end{itemize}
Now we can state that for $f=h+g$, then
\[
\text{LHS}\to Pf = h \text{ as } n\to\infty
\]

\subsection{Proving the orthogonal complement}

What we didn't show was how $H^{\text{inv}}= G^\bot$.

\textbf{First direction:} Let $g=Uf-f$. If $h\in H^{\text{inv}}$, then $h\bot g$ for all $g\in G$.
\begin{align*}
    \langle h, Uf-f\rangle &= \langle h, Uf \rangle - \langle h, f\rangle \\
    &= \langle Uh, Uf \rangle - \langle h, f\rangle \quad h\text{ invariant} \\
    &= \langle h, f \rangle - \langle h, f\rangle \quad U\text{ isometry} \\
    &= 0 \implies h\in G^\bot
\end{align*}

\textbf{Second direction:} Take $h\in G^\bot$. We need to show $Uh=h$. For all $f\in H$, it follows that 
\begin{align*}
    &\langle h, Uf-f\rangle=0\\
    \implies& \langle h, Uf\rangle = \langle h, f \rangle \quad \forall f\in H\\
    \implies& \langle U^\ast h, f\rangle = \langle h, f \rangle \quad \forall f\in H\\
    \implies& \langle U^\ast h - h, f \rangle =0\\
    \implies& U^\ast h = h \implies Uh=h
\end{align*}
We need a small lemma to show that $U^\ast h = h \implies Uh=h$.

\subsubsection{Lemma}

Suppose $U$ is an isometry. Then 
\[
Uf =f \iff U^\ast f = f
\]
For the first direction, we assume $Uf = f$ and then we write 
\[
U^\ast Uf = U^\ast f \implies f = U^\ast f
\]
where $U^\ast U=I$. For the second direction, we assume $U^\ast = f$ and then if we can prove that 
\[
\left|\left| Uf-f \right|\right|^2 = 0
\]
where it follows that 
\[
Uf=f
\]
Hence, we write 
\begin{align*}
    \langle Uf-f, Uf-f \rangle &= \langle Uf, Uf \rangle - \langle f, Uf \rangle -\langle Uf, f \rangle + \langle f, f \rangle\\
    &= \langle f, f \rangle - \langle U^\ast f, f \rangle - \langle f, U^\ast f \rangle + \langle f, f \rangle\\
    &= \langle f, f \rangle - \langle f, f \rangle - \langle f, f \rangle + \langle f, f \rangle=0
\end{align*}

\section{Birkhoff's ergodic theorem}

Birkhoff's (``pointwise'') ergodic theorem states that for 
\[
\lim_{n\to\infty} \frac{1}{n}\left(f(x)+f(Tx)+\cdots + f(T^{n-1}x)\right)\triangleq \bar{f}(x)
\]
Also, 
\[
\bar{f}(Tx) = \bar{f}(x)
\]
Then, $\bar{f}\in L'(M)$ and 
\[
\int \bar{f}\, d\mu = \int f\, d\mu
\]
which is the spatial average of $f$.

\section{Invariant sets and functions}

A function $f$ is invariant if $f(Tx)=f(x)$ for all $x\in M$ where $f$ is an endomorphism and $T$ is an automorphism
\[
f(Tx) = f(x) = f(T^{-1}x)
\]
This just means that $f$ is a constant along a trajectory.

Set $A$ is invariant if $\mathbbm{1}_A$ is an invariant function
\[
\mathbbm{1}_A(x)=\begin{cases}
    1 \text{ if } x\in A\\
    0 \text{ else}
\end{cases}
\]
An example of this would be the unit circle 
\[
Tx = x+\alpha \pmod 1
\]
If $\alpha$ is irrational, then the only invariant sets at $\emptyset$ or $M$.

\section{What does it mean to be ergodic}

\underline{Definition}: $T$ is ergodic if the only invariant sets have full measure, or zero measure
\[
\mu(A) = 1 \quad \text{or} \quad \mu(A)=0
\]
\underline{Theorem}: If $T$ is ergodic, then any invariant function is constant on a set of full measure.

Proof: Suppose $T$ is ergodic. Then $Pf$ equals to the projection onto constant functions
\[
Pf = \frac{\langle f, 1\rangle}{\langle 1, 1\rangle}\cdot 1 = \int f(x)\cdot 1 \, d\mu(x) = \int f\, d\mu
\]
where $\int f\, d\mu$ is the spatial average.

\subsection{A note}

At this point, Rowley made the explicit statement that 
\[
\lim_{n\to\infty}\,\frac{1}{n}\left(f+ Uf + U^2 f +\ldots + U^{n-1}f\right)
\]
is the time average and 
\[ Pf\]
is the spatial average. Ergodic theorems are about relating time averages to spatial averages.

\subsection{Equivalent notions of ergodicity}
\begin{itemize}
    \item Each invariant function ($f(x)=f(Tx)$ almost everywhere) is constant almost everywhere
    \item In the Birkhoff theorem, the limiting $\bar{f}(x)$ is constant. In other words \[ \frac{1}{n}\left(f(x) + f(Tx) +\cdots + f(T^{n-1}x)=\bar{f}=\int fd\mu\right) \] where the left hand side is the time average and the right hand side is the spatial average.
\end{itemize}

\section{Back to our example of rotations}

Suppose $\alpha$ is irrational in 
\[
Tx = x+\alpha \pmod 1
\]
Let $f$ be an invariant function. Write 
\[
f=\sum_{k=-\infty}^{\infty} f_k e^{2\pi i k x}
\]
and then 
\begin{align*}
    Uf(x) &= \sum_{k=-\infty}^{\infty}f_k e^{2\pi i k (x+\alpha)}\quad \text{where invariant functions $Uf=f$}\\
    &= \sum_{k=-\infty}^{\infty}f_k e^{2\pi i k x}\implies \sum_k f_k \left(e^{2\pi i k x}-e^{2\pi i k (x+\alpha)}\right)=0\\
    &\implies \sum_k f_k e^{2\pi i k x}\left( 1-e^{2\pi i k \alpha} \right)=0
\end{align*}
Hence 
\[
\sum_k f_k \int_{0}^{2\pi}e^{2\pi i (k-m)x}\, dx\, \left(1-e^{2\pi i k \alpha}\right)\implies f_m \left(1-e^{2\pi i m \alpha}\right) = 0
\]
Say that $f_m \left(1-e^{2\pi i m \alpha}\right) = 0$ because $1-e^{2\pi i m \alpha}=0$. This implies
\[
e^{2\pi i m \alpha} = e^{2\pi i n}
\]
where it follows that $\alpha=\frac{n}{m}$, violating our assumption. Hence, we have
\[
f_m=0 \quad \forall\, m\neq 0 
\]
where we note that $f_0$ can be nonzero. Therefore, we have 
\[
f(x)=f_0
\]
where $f_0$ is a constant function.

\section{How can we state if our ergodic systems are chaotic?}

We need to look at the properties of the eigenvalues and eigenfunctions of $U$.

Let's start by writing 
\[
Uf(x) = f(Tx)
\]
where $T$ is measure preserving, which implies $\left|\left| Uf \right|\right| = \left|\left| f \right|\right|$. Now we are going to make 4 claims that will allow us to talk about the properties of ergodic functions in terms of their eigenfunctions.

\subsection{Claim 1}

The eigenvalues of $U$ must be on the unit circle. For a proof, suppose $Uf=\lambda f$ for $f\neq 0$ and $\lambda\in\mathbb{C}$. Then we have 
\begin{gather*}
    \langle Uf, Uf\rangle = \langle \lambda f, \lambda f \rangle\\
    \langle f, f\rangle = \lambda\bar{\lambda} \langle f, f\rangle \\
    \implies \lambda\bar{\lambda} = 1 \text{ and }\left|\lambda\right|^2 =1
\end{gather*}
This holds for an isometry over a complex field.

\subsection{Claim 2}

For constant functions such as $f\equiv c$, we have 
\[
Uf(x)=f(Tx) = c = f(x)
\]
Therefore
\[
Uf = f
\]
which implies $f$ is an eigenfunction of $U$ with eigenvalue $1$.

\subsection{Claim 3}

If $T$ is ergodic, then constant functions are the only eigenfunctions with eigenvalue of 1. This follows from the definition of ergodicity. \textbf{Ergodicity is literally the statement that the eigenvalue 1 is simple.} 

From this claim, we will begin to think about different ergodic systems $T$ in terms of the spectral properties of $U_T$.

\subsection{Claim 4}

If $Uf=e^{i\theta}f$ and $\theta\neq 0$, then 
\[
\int f\, d\mu=0.
\]

For a proof, note that eigenvectors (and hence, eigenfunctions) of an isometry for distinct eigenvalues are orthogonal. From \text{claim 2}, we know that $1$ is an eigenfunction of $U_T$; therefore, $f\bot 1$. Hence we can write 
\[
0=\langle f, 1 \rangle = \int f\cdot \bar{1}\, d\mu = \int f\, d\mu
\]

For an alternative proof, we can write 
\[
\int f \, d\mu = \int f \circ T\, d\mu = \int Uf \, d\mu = \int e^{i\theta}f\, d\mu = e^{i\theta}\int f \, d\mu
\]
which boils down to 
\[
\int f \, d\mu = e^{i\theta}\int f \, d\mu
\]
which implies 
\[
\left( 1- e^{i\theta} \right)\int f \, d\mu = 0
\]
Because $\theta\neq 0$, it follows that $\left( 1- e^{i\theta} \right)\neq 0$, implying 
\[
\int f \, d\mu = 0
\]
This gives us license to only look at functions in $H_0$ defined as
\[
H_0 = \left\{f : \int f \, d\mu = 0 \right\}
\]
i.e. the non-constant functions, when studying different ergodic systems.

\section{Mixing properties}

Because a system is Ergodic, we can write
\[
\lim_{n\to\infty} \frac{1}{n}\sum_{k=0}^{n-1}f(T^k x)=\int_M f\, d\mu
\]
If we multiply by $g\in L^2$ and integrate, we have 
\[
\lim_{n\to\infty} \frac{1}{n}\sum_{k=0}^{n-1} \int_M f(T^k x) g(x) \,d\mu(x) = \int_M f\,d\mu \cdot \int g\, d\mu
\]

The definition of weak mixing for all $f,g$ is 
\[
\lim_{n\to \infty}\frac{1}{n}\sum_{k=0}^{n-1}\left| \int_{M} f\left(T^k x\right)g(x)\,d\mu - \int_M f\,d\mu\,\cdot\, \int g\,d\mu \right|=0
\]

The definition of mixing is for all $f,g$ we have
\[
\lim_{n\to \infty}\int_M f\left(T^n x\right) g(x)\, d\mu(x) = \int f\,d\mu \,\cdot\, \int g\, d\mu 
\]
where we now have the following relationship
\[
\text{Mixing}\implies \text{weak mixing} \implies ergodic
\]

\subsection{Why mixing}

Let $f=\mathbbm{1}_A$ and $g=\mathbbm{1}_B$. Then mixing implies 
\[
\lim_{n\to \infty }\mu\left(T^{-n}A\cap B\right)=\mu(A)\cdot\mu(B)
\]

What's interesting is that for the map
\[
Tx = x+\alpha \pmod 1
\]
when $\alpha$ is rational, the system is not ergodic, but when $\alpha$ is irrational, the system is ergodic but is not mixing.

\subsection{Doubling map}

The doubling map 
\[
Tx = 2x \pmod 1
\]
is mixing, hence, is ergodic.

\subsubsection{A theorem}

If $Uf=e^{i\theta}f$ where $\theta\neq 0$ and $f\neq 0$, then $T$ cannot be weak mixing.

For a proof, we can write 
\begin{align*}
    \left|\langle U^k f, g\rangle - \int f\,d\mu\,\cdot\, \int g\,d\mu\right| &= \left|\langle U^k f, g\rangle\right| = \left|\langle e^{ik\theta} f, g\rangle\right|\\
    &= |\langle f, g \rangle| \neq 0 \quad (\text{take }g=f)
\end{align*}
where $\int f\,d\mu\,\cdot\, \int g\,d\mu=0$. Hence mixing is a spectral property of $U$.



\end{document}